\section{Week in Review}
\sectionkicker{WEEKLY SYNTHESIS}
\begin{wideflow}
{\Large\bfseries 今週は、モデルの利用条件まで動いた}\par
\smallskip
{\color{SurveyMuted}モデル/APIの入口、runtime、agent・multimodalの測り方までが一つの流れとして見えてきた。使えること、動かせること、確かめられることを横断して読む。}\par
\end{wideflow}

\subsection*{今週何が変わったか}
\begin{enumerate}[leftmargin=1.6em]
\item \textbf{入口が分岐した。} API preview、GA API/app/web、open weights、partner channelが並び、モデルの存在と利用可能性を同じ語で扱えなくなった。Cyberのprogram accessも含め、誰がどの目的でどの経路を通るかがavailabilityの一部になった。
\item \textbf{運用の層が動いた。} vLLM、llama.cpp、SGLang、FlashInferの更新は、serving framework、local runtime、front-end/cache、low-level kernelが別々の実装面で進むことを示した。KVの配置・prefetchやdecoding policyまでが「モデルを動かす」条件として前景化した。
\item \textbf{確かめる単位が細かくなった。} agentではscaffolding、requirements/planning、function call、transaction、red teaming、skillの失敗箇所を分け、multimodalでは理解・生成編集・runtimeの接点を分けて測る方向が見えた。
\end{enumerate}

\subsection*{なぜそれらが一緒に重要なのか}
新しいaccess surfaceは、対応するruntimeがあって初めて運用へ届く。runtimeの効率は、KVの仮想化、tiered memory、sparse prefetch、decodingの選択に左右される。agentやmultimodalの実用性は、tool、状態、workflowを含む評価で初めて確かめられる。今週の変化は、モデルを手に入れる入口、動かし方、壊れ方の測り方が別々の話ではなく、利用条件を構成する連鎖として見えるようになった点にある。

この連鎖は製品間の直接的な互換性を意味しない。各release、paper、program recordは固有の主体・環境・測定条件を持つ。accessの拡大を性能の向上へ、runtimeの更新を全モデルの対応へ、細かなbenchmarkを一般のagent reliabilityへ自動的に変換してはいけない。

\begin{communitynote}[Weekly Community Movement --- context only]
X上の話題では、GLM-5.3、Grok 4.6、Qwen3.8への関心と、ローカル推論、coding/agent、価格競争への関心が並んで見られた。これは今週の関心の広がりを示す背景であり、モデルの性能、提供条件、互換性を証明するものではない\cite{w33-community}。
\end{communitynote}

\subsection*{次に何を見るべきか}
次週以降は、この連鎖のどこがさらに具体化するかを追う。UltrafastのpreviewとGA、速度値、GLM-5.3のdirect-page・benchmark・cyber・local-weight timing、Daybreakのmodel/access scopeとBedrock境界はまだ詳細化されていない。VoiceDesignerもmodel/data contribution、baseline、evaluation、noveltyの捕捉が足りない。さらに、runtimeや論文の評価が、それぞれのモデル・環境・条件から外へ一般化できるかを確認する必要がある。

\begin{claimboundary}[この総括の読み方]
各項目の公開条件、報告主体、測定環境は同じではない。提供元が示す速度、論文著者が報告する評価、コミュニティで見られた関心は、それぞれ異なる種類の情報であり、単純な性能ランキングへまとめられない。未確定の点は、次の公開更新や再現可能な評価で確かめるべき観測点として残る。
\end{claimboundary}